\documentclass[a4paper,11pt,oneside]{memoir}

% Castellano
\usepackage[spanish]{babel}
\selectlanguage{spanish}
\usepackage[utf8]{inputenc}
\usepackage{placeins}
\usepackage[gen]{eurosym}
\usepackage{lscape}

\RequirePackage{booktabs}
\RequirePackage[table]{xcolor}
\RequirePackage{xtab}
\RequirePackage{multirow}

% Links
\usepackage[colorlinks]{hyperref}
\hypersetup{
	allcolors = {red}
}

% Ecuaciones
\usepackage{amsmath}

% Rutas de fichero / paquete
\newcommand{\ruta}[1]{{\sffamily #1}}

% Párrafos
\nonzeroparskip

% Profundidad de índice
\setcounter{tocdepth}{2}


% Imagenes
\usepackage{graphicx}
\newcommand{\imagen}[2]{
	\begin{figure}[!h]
		\centering
		\includegraphics[width=0.9\textwidth]{#1}
		\caption{#2}\label{fig:#1}
	\end{figure}
	\FloatBarrier
}

\newcommand{\imagenflotante}[2]{
	\begin{figure}%[!h]
		\centering
		\includegraphics[width=0.9\textwidth]{#1}
		\caption{#2}\label{fig:#1}
	\end{figure}
}



% El comando \figura nos permite insertar figuras comodamente, y utilizando
% siempre el mismo formato. Los parametros son:
% 1 -> Porcentaje del ancho de página que ocupará la figura (de 0 a 1)
% 2 --> Fichero de la imagen
% 3 --> Texto a pie de imagen
% 4 --> Etiqueta (label) para referencias
% 5 --> Opciones que queramos pasarle al \includegraphics
% 6 --> Opciones de posicionamiento a pasarle a \begin{figure}
\newcommand{\figuraConPosicion}[6]{%
  \setlength{\anchoFloat}{#1\textwidth}%
  \addtolength{\anchoFloat}{-4\fboxsep}%
  \setlength{\anchoFigura}{\anchoFloat}%
  \begin{figure}[#6]
    \begin{center}%
      \Ovalbox{%
        \begin{minipage}{\anchoFloat}%
          \begin{center}%
            \includegraphics[width=\anchoFigura,#5]{#2}%
            \caption{#3}%
            \label{#4}%
          \end{center}%
        \end{minipage}
      }%
    \end{center}%
  \end{figure}%
}

%
% Comando para incluir imágenes en formato apaisado (sin marco).
\newcommand{\figuraApaisadaSinMarco}[5]{%
  \begin{figure}%
    \begin{center}%
    \includegraphics[angle=90,height=#1\textheight,#5]{#2}%
    \caption{#3}%
    \label{#4}%
    \end{center}%
  \end{figure}%
}
% Para las tablas
\newcommand{\otoprule}{\midrule [\heavyrulewidth]}
%
% Nuevo comando para tablas pequeñas (menos de una página).
\newcommand{\tablaSmall}[5]{%
 \begin{table}
  \begin{center}
   \rowcolors {2}{gray!35}{}
   \begin{tabular}{#2}
    \toprule
    #4
    \otoprule
    #5
    \bottomrule
   \end{tabular}
   \caption{#1}
   \label{tabla:#3}
  \end{center}
 \end{table}
}

%
% Nuevo comando para tablas pequeñas (menos de una página).
\newcommand{\tablaSmallSinColores}[5]{%
 \begin{table}[H]
  \begin{center}
   \begin{tabular}{#2}
    \toprule
    #4
    \otoprule
    #5
    \bottomrule
   \end{tabular}
   \caption{#1}
   \label{tabla:#3}
  \end{center}
 \end{table}
}

\newcommand{\tablaApaisadaSmall}[5]{%
\begin{landscape}
  \begin{table}
   \begin{center}
    \rowcolors {2}{gray!35}{}
    \begin{tabular}{#2}
     \toprule
     #4
     \otoprule
     #5
     \bottomrule
    \end{tabular}
    \caption{#1}
    \label{tabla:#3}
   \end{center}
  \end{table}
\end{landscape}
}

%
% Nuevo comando para tablas grandes con cabecera y filas alternas coloreadas en gris.
\newcommand{\tabla}[6]{%
  \begin{center}
    \tablefirsthead{
      \toprule
      #5
      \otoprule
    }
    \tablehead{
      \multicolumn{#3}{l}{\small\sl continúa desde la página anterior}\\
      \toprule
      #5
      \otoprule
    }
    \tabletail{
      \hline
      \multicolumn{#3}{r}{\small\sl continúa en la página siguiente}\\
    }
    \tablelasttail{
      \hline
    }
    \bottomcaption{#1}
    \rowcolors {2}{gray!35}{}
    \begin{xtabular}{#2}
      #6
      \bottomrule
    \end{xtabular}
    \label{tabla:#4}
  \end{center}
}

%
% Nuevo comando para tablas grandes con cabecera.
\newcommand{\tablaSinColores}[6]{%
  \begin{center}
    \tablefirsthead{
      \toprule
      #5
      \otoprule
    }
    \tablehead{
      \multicolumn{#3}{l}{\small\sl continúa desde la página anterior}\\
      \toprule
      #5
      \otoprule
    }
    \tabletail{
      \hline
      \multicolumn{#3}{r}{\small\sl continúa en la página siguiente}\\
    }
    \tablelasttail{
      \hline
    }
    \bottomcaption{#1}
    \begin{xtabular}{#2}
      #6
      \bottomrule
    \end{xtabular}
    \label{tabla:#4}
  \end{center}
}

%
% Nuevo comando para tablas grandes sin cabecera.
\newcommand{\tablaSinCabecera}[5]{%
  \begin{center}
    \tablefirsthead{
      \toprule
    }
    \tablehead{
      \multicolumn{#3}{l}{\small\sl continúa desde la página anterior}\\
      \hline
    }
    \tabletail{
      \hline
      \multicolumn{#3}{r}{\small\sl continúa en la página siguiente}\\
    }
    \tablelasttail{
      \hline
    }
    \bottomcaption{#1}
  \begin{xtabular}{#2}
    #5
   \bottomrule
  \end{xtabular}
  \label{tabla:#4}
  \end{center}
}



\definecolor{cgoLight}{HTML}{EEEEEE}
\definecolor{cgoExtralight}{HTML}{FFFFFF}

%
% Nuevo comando para tablas grandes sin cabecera.
\newcommand{\tablaSinCabeceraConBandas}[5]{%
  \begin{center}
    \tablefirsthead{
      \toprule
    }
    \tablehead{
      \multicolumn{#3}{l}{\small\sl continúa desde la página anterior}\\
      \hline
    }
    \tabletail{
      \hline
      \multicolumn{#3}{r}{\small\sl continúa en la página siguiente}\\
    }
    \tablelasttail{
      \hline
    }
    \bottomcaption{#1}
    \rowcolors[]{1}{cgoExtralight}{cgoLight}

  \begin{xtabular}{#2}
    #5
   \bottomrule
  \end{xtabular}
  \label{tabla:#4}
  \end{center}
}




\graphicspath{ {./img/} }

% Capítulos
\chapterstyle{bianchi}
\newcommand{\capitulo}[2]{
	\setcounter{chapter}{#1}
	\setcounter{section}{0}
	\chapter*{#2}
	\addcontentsline{toc}{chapter}{#2}
	\markboth{#2}{#2}
}

% Apéndices
\renewcommand{\appendixname}{Apéndice}
\renewcommand*\cftappendixname{\appendixname}

\newcommand{\apendice}[1]{
	%\renewcommand{\thechapter}{A}
	\chapter{#1}
}

\renewcommand*\cftappendixname{\appendixname\ }

% Formato de portada
\makeatletter
\usepackage{xcolor}
\newcommand{\tutor}[1]{\def\@tutor{#1}}
\newcommand{\course}[1]{\def\@course{#1}}
\definecolor{cpardoBox}{HTML}{E6E6FF}
\def\maketitle{
  \null
  \thispagestyle{empty}
  % Cabecera ----------------
\noindent\includegraphics[width=\textwidth]{cabecera}\vspace{1cm}%
  \vfill
  % Título proyecto y escudo informática ----------------
  \colorbox{cpardoBox}{%
    \begin{minipage}{.8\textwidth}
      \vspace{.5cm}\Large
      \begin{center}
      \textbf{TFG del Grado en Ingeniería Informática}\vspace{.6cm}\\
      \textbf{\LARGE\@title{}}
      \end{center}
      \vspace{.2cm}
    \end{minipage}

  }%
  \hfill\begin{minipage}{.20\textwidth}
    \includegraphics[width=\textwidth]{escudoInfor}
  \end{minipage}
  \vfill
  % Datos de alumno, curso y tutores ------------------
  \begin{center}%
  {%
    \noindent\LARGE
    Presentado por \@author{}\\ 
    en Universidad de Burgos --- \@date{}\\
    Tutor: \@tutor{}\\
  }%
  \end{center}%
  \null
  \cleardoublepage
  }
\makeatother


% Datos de portada
\title{Aplicación web para la visualización y búsqueda de datos almacenados en dispositivo con linux embebido.}
\author{Rodrigo Martínez Bravo}
\tutor{Jesús Enrique Sierra García}
\date{\today}

\begin{document}

\maketitle



\cleardoublepage



%%%%%%%%%%%%%%%%%%%%%%%%%%%%%%%%%%%%%%%%%%%%%%%%%%%%%%%%%%%%%%%%%%%%%%%%%%%%%%%%%%%%%%%%



\frontmatter


\clearpage

% Indices
\tableofcontents

\clearpage

\listoffigures

\clearpage

%\listoftables

%\clearpage

\mainmatter

\appendix

\apendice{Plan de Proyecto Software}

\section{Introducción}
En este apartado se especificará la planificación seguida en cuanto a tiempo se refiere para realizar cada una de las tareas que han conformado este proyecto, permitiendo gracias a esta planificación cumplir con los objetivos que se han establecido para el sistema desarrollado.\\
Debido a que gran parte de este proyecto se ha basado en la investigación (ya que, como se ha explicado en otros puntos, se desconocían algunos de los procedimientos a realizar), la planificación temporal no ha podido ser tan estricta como debería ser en un proyecto cuya organización se basa en la metodología ágil conocida como SCRUM como es este caso. Otro aspecto que ha impedido seguir esa planificación ha sido la aparición de una gran cantidad de contratiempos surgidos durante cada una de las fases de este proyecto, como se describirá más adelante.\\
En cuanto al estudio de viabilidad, la parte económica ha sido también un aspecto difícil de calcular ya que se ha necesitado probar una serie de implementaciones, tanto hardware como software, hasta dar con la idónea para cumplir los requisitos del proyecto.\\
 
\section{Planificación temporal}
Como bien se ha descrito en la introducción, en este apartado se especificará el desarrollo del proyecto descrito a partir de los periodos de tiempo que han llevado a desarrollar cada parte, o como se conoce en terminología SCRUM, cada uno de los sprints.\\
\subsection{Sprint 0: 05/10/2017 - 22/11/2017}
Este es el sprint inicial, donde se explica y establece los conceptos sobre los que se va a trabajar y los elementos que se van a emplear para ello. Se plante la realización de este trabajo basandose en el sistema operativo Ubuntu, en su versión 14.04 LTS, para desarrollar el código usando la aplicación conocida como Eclipse en su versión 4.8 (conocida como Photon) y el leguaje de programación C$++$. Además se plantearon las librerías a utilizar para este proceso y la utilización del dispositivo MTX-GTW para la ejecución remota del sistema en este dispositivo.\\
\subsubsection{Tareas}
\begin{itemize}
	\item Creación de la máquina virtual con el sistema operativo antes mencionado.
	\item Configuración de Eclipse con las correspondientes librerías.
	\item Preparación de los diagramas que servirán como guía para el desarrollo del sistema.
	\item Búsqueda de cables de alimentación para cada uno de los elementos.
	\item Comienzo de investigación sobre la realización de ejecuciones remotas.
	\item Creación del repositorio.
\end{itemize}
\subsubsection{Backlog}
En este sprint inicial se plantea cada uno de los aspectos del proyecto (requisitos, objetivos...) así como el planteamiento de la organización con el tutor en cuanto a reuniones que se iban a realizar. Además se empezó a organizar la estructura del programa con la realización de los diagramas de clases.\\
\subsubsection{Investigación}
Se comienza la investigación a cerca de las ejecuciones remotas de programas para poder realizarlo con el MTX-GTW. Tambien se comienza a investigar a cerca del manejo de sockets y lecturas de puertos utilizando C$++$ para poder recoger los datos del láser y enviar mensajes por este mismo puerto para comunicarse con el láser (con la dificultad de que esta escucha también será remota).\\

\subsection{Sprint 1: 22/11/2017 - 16/01/2018}
En este sprint se plantean los cambios a realizar en el diagrama creado en el anterior sprint, se plantea el uso de \LaTeX para la realización de la memoria del proyecto y la creación del diagrama de secuencia para mejor comprensión del funcionamiento del sistema. Se comprueba también el estado de la máquina virtual y se detecta la falta de algunas librerías necesarias. \\
\subsubsection{Tareas}
\begin{itemize}
	\item Creación del diagrama de secuencia.
	\item Mejora del diagrama de clases según la revisión realizada.
	\item Solucionar la ausencia de librerías antes mencionado.
	\item Realización de pruebas de conexión con el MTX-GTW.
	\item Continuación de búsqueda de material hardware.
\end{itemize}
\subsubsection{Backlog}
En este sprint se termina de crear y modificar los diagramas que servirán de base a la posterior creación del código del sistema, se pretende configurar Eclipse con las librerías que faltaban del sprint anterior y se implementa el programa que imprima por pantalla "Hola mundo" realizando esta ejecución en el dispositivo remoto. Mientras se realizan estos procesos se prosigue con la búsqueda del cable de alimentación del láser ya que el cable del MTX-GTW ya ha sido encontrado.\\
\subsubsection{Investigación}
En este sprint se realizan varias tareas de investigación. En primer lugar, se deben encontrar las librerías que se necesitan para incluirlas en la configuración del proyecto en Eclipse. Esta tarea resulta costosa debido a que algunas de las librerías que se habían acordado utilizar estaban desactualizadas y no se podían encontrar y se desconocía la organización de ficheros y el funcionamiento de las nuevas versiones con lo que se hacia complicado la configuración para algunas de las tareas futuras.\\
La otra tarea de investigación realizada fue la búsqueda de los diferentes cables de alimentación y esto es debido a la diferencia entre ambos ya que el dispositivo MTX-GTW necesita ser alimentado con una corriente alterna de 12 V. Debido a que, en la actualidad, la mayoría de los cargadores de los diferentes aparatos electrónicos cotidianos no suministran más de 7 voltios se necesitó fabricar un cargador a base de un transformador que suministrase la corriente deseada unido a dos cable los cuales irán conectados al cabezal del aparato debido a que tampoco posee un cabezal convencional. Por otro lado, el láser  se han utilizado dos cableados distintos ya que se debe conectar tanto a la red eléctrica para alimentarse como al dispositivo antes mencionado. Para la primera de las conexiones descritas se ha empleado un adaptador a corriente alterna ya que este láser necesita alimentarse con 24 V de corriente alterna. Para la segunda conexión se necesita un cable con cabezal VGA y el otro extremos sin cabezal para poder conectar solo los cables que sean necesarios ya que solo se necesitan algunos de los pines de esta conexión para el envío y recepción de datos.\\

\subsubsection{Pruebas}
En este sprint se comienzan a realizar las primeras pruebas para la conexión del PC con el MTW-GTW para configurar de forma correcta esta conexión con el fin de tenerla de forma definitiva como conexión de ambos sistemas en el proyecto final.\\

\subsection{Sprint 2: 16/01/2018 - 01/02/2018}
Una vez se han creado los diagramas para aclarar en funcionamiento del sistema final, se plantean las tareas para este nuevo sprint. En este caso, una vez se posee una versión estable del programa que realiza las pruebas se empieza a plantear el inicio del desarrollo del código del sistema, así como la creación de una cuenta en la web Trello (mencionada en puntos anteriores) para la creación del panel con las tareas a realizar con lo que se consigue una version del panel de tareas para la organización de cada sprint.\\
\subsubsection{Tareas}
\begin{itemize}
	\item Paso del proyecto de pruebas al repositorio para comenzar en el el desarrollo del sistema 
	\item Creación de cuenta y tablón en Trello.
	\item Creacion de las primeras clases.
\end{itemize}
\subsubsection{Backlog}
En este sprint se ha comenzado con el desarrollo del sistema ya en el repositorio para mantener un control de las versiones que se van realizando. Tambien se comienza a dejar constancia de cada una de las tareas que se han plantado en los sprint, dividiéndolas en diferentes categorias segun se encuentren pendientes de realizar, en curso o si ya han sido realizadas.
\subsubsection{Investigación}
Se retoma la investigación a cerca de los sockets en el lenguaje de C$++$ para plantear la conexión con el láser con el objetivo de organizar la estructura final de la comunicación lógica de los diferentes elementos físicos que componen el sistema.

\subsection{Sprint 3: 01/02/2018 - 20/05/2018}
Este es uno de los sprints más largos debido a que ha sido en el que han surgido la mayoria de problemas. En este punto, al no poder plantear nuevas tareas debido al bloqueo de las ya presentes del anterior sprint, las reuniones se fueron separando unas de otras en el tiempo. 
\subsubsection{Tareas}
\begin{itemize}
	\item Solución de los problemas de Ubuntu y Eclipse.
	\item Continuación de la creación de las distintas clases y procedimientos.
\end{itemize}
\subsubsection{Backlog}
En este Sprint, como se ha dicho anteriormente, surgen la mayor parte de las incidencias que impiden el avance ene el proyecto por un largo periodo de tiempo. Una vez en proceso de creación del sistema, Eclipse comenzó a dar errores en cuanto a la comprensión de comandos básicos del lenguaje C$++$. Junto con esta clase de problemas surgió el inconveniente del fallo completo de Eclipse, ya que en la máquina virtual donde se estaba realizando todo el proceso, la aplicación dejó de funcionar, permitiendo verse en pantalla únicamente la barra de tareas. Para solucionar este fallo se realizaron algunas investigaciones sin éxito.\\
\subsubsection{Investigación}
Las investigaciones realizadas en este periodo de tiempo fueron varias aunque todas con el mismo objetivo, solucionar fallos.\\
En un inicio, los fallos en cuanto al lenguaje de programación. Para solucionar este tipo de fallo se realizón un análisis exaustivo de la configuración completa  del proyecto, incluyendo partes que no se habían visto hasta ese momento, por si se podría haber producido algún efecto en cascada que hubiese desencadenado el error (dicho análisis fue realizado de forma conjunta con el profesor). Tras comprobar que ese no era el problema se comenzó una búsqueda por diversas páginas web para encontrar soluciones de otros usuarios que hubiesen tenido el mismo problema. Al realizar esta búsqueda, no solo se descubrió que el número de personas cone este error era muy reducido, sino que además se pudo observar que las soluciones que se les ofrecía para solucionarlo o no funcionaban o no eran útiles para el problema que se planteaba en nuestro caso. Como conclusión a este problema se decidió crear otro proyecto y dejar para más adelante la configuración realizada con anterioridad.\\
En cuanto al problema de la visualización de Eclipse, la situación era similar a la del anterior problema. En este caso el número de personas con el mismo problema era mayor pero la solución que se daba en la mayoría de los casos era la eliminación de los recursos temporales generados al arrancar Eclipse para reiniciarlo y que los crease de nuevo, intentando solventar así los fallos que hayan podido aparecer durante el arranque. Al no conseguir solucionarlo de esta manera, se optó por la opción más drástica, extraer los cambios ya realizados en el proyecto creado para evitar el primer problema, eliminar la máquina virtual y comenzar la instalación y configuración de aplicaciones desde el principio.\\
Este método no funcionó, por lo que se decidió tomar un camino alternativo para el proyecto.

\subsection{Sprint 4: 20/05/2018 - 13/06/2018}
Tras el periodo anterior de cambios de máquina virtual y reconfiguraciones se tomó la decisión de cambiar el proyecto. A partir de este sprint el proyecto sería realizado en Windows 10 (sistema anfitrión del PC) para evitar problemas con los servicios de virtualización. Debido a esa decisión se hacia incompatible la utilización del MTX-GTW, ya que este aparato es incompatible con Windows, por lo que la ejecución del sistema se realizaría en el PC. Este aspecto solucionó otro de los problemas anteriores, las librerías. Al no tener que realizar ninguna ejecución remota, el proyecto solo iba a necesitar utilizar las librerías que ya vienen por defecto en C$++$, sin necesidad de descargar ninguna otra. En este periodo se comienza las primeras interacciones con el láser
\subsubsection{Tareas}
\begin{itemize}
	\item Configuración del proyecto y desarrollo del código en Windows.
	\item Investigación sobre programas de manejo del láser.
	\item Cambios importantes en la memoria.
\end{itemize}
\subsubsection{Backlog}
En este periodo de tiempo se consigue la realización de gran parte del código que forma el sistema software ya que no se producen errores como los del anterior sprint. También se instala y empieza a investigar el funcionamiento del programa que se encuentra adjunto al láser en un CD presente en la misma caja del aparato. Junto con este programa, tambien se descarga e instala el programa RealTerm, el cual se usará para envir y recibir mensajes por el puerto USB, ya que el láser tiene la posibilidad de conectarse al PC por esta vía.
\subsubsection{Investigación}
En este aspecto se empieza a investigar la configuración del programa RealTerm antes mencionado, dejando para más adelante la programación de la parte restante del sistema, para centrarse en averiguar como conseguir la conexión con el láser además de conoces como realizar el intercambio de información.\\
Tambien se investiga la forma de configurar del programa del láser (UAM Proyect Designer) para poder ver el funcionamiento de este programa y la visualización por pantalla de las lecturas del láser, lo cual se podrá tomar como guía en la representación del resultado de la ejecución del sistema objetivo. Esto se consigue gracias a la configuración de los puertos USB de PC para que permitan instalar drivers no firmados (como son los del propio láser, los cuales tambien se tuvieron que buscar por internet para evitar la dependencia con el CD) con lo que se pudiese reconocer al láser para que el programa funcionase correctamente. Este buen comportamiento también se produjo ya que el tutor facilitó un fichero de configuración básica el cual fue enviado vía USB al propio láser, el cual lo conserva en memoria desde entonces.
\subsection{Sprint 5: 13/06/2018 - 19/11/2018}
En este sprint, cuya duración se debe a la aparición del último de los problemas de gran relevancia para el proyecto y la presencia del periodo de verano lo cual supone una mayor dificultad para las reuniones. En este sprint se analizan los resultados del ciclo anterior, observando las lecturas del láser por primera vez desde el inicio del proyecto. El problema antes mencionado surge al intentar crear el socket lógico en C$++$ tras conocer como crear los mensajes y provocar la comunicación entre los dos aparatos y conseguir asi implementar una de las partes más importantes del proyecto. No se consigue (entre el alumno y el profesor) que el socket conecte el PC con el láser con lo que se dedica este periodo a solucionar este problema
\subsubsection{Tareas}
\begin{itemize}
	\item Solucionar falta de comunicación con el láser desde el PC.
	\item Cambio de lenguaje del proyecto
\end{itemize}
\subsubsection{Backlog}
En este periodo se realizan algunas investigaciones con resultados nada satisfactorios y, tras no conseguir el bojetivo principal de este sprint, se decide cambiar de lenguaje. En este caso se pasa todo el código existente a Python. Este lenguaje, al ser más usado, posee un soporte mucho mayor, especialmente por parte de los usuarios que lo usan para desarrollar su código, por lo que la traducción del código se realiza de forma rápida y sencilla. Gracias a la utilización de este lenguaje se consigue configurar el socket y tenerlo plenamente operativo de forma muy sencilla, con lo que en el periodo final de este sprint se consigue solucionar el priblema inicial y provocar un grán avance.
\subsubsection{Investigación}
En este periodo la primera investigación, como ya se ha comentado anteriormente, es la solución del socket en el lenguaje de C$++$, sin resultado satisfactorio. A parte de esta, la investigación principal se centró en la creación de sockets en Python lo cual se consiguió de forma muy sencilla gracias a la ayuda proporcionada por la propia web del lenguaje y los usuarios que poseian este problema en diferentes foros web.
\subsubsection{Pruebas}
Se realizan las primeras pruebas de comunicación y obtención de respuestas por parte del láser que da como resultado la implementación del procedimiento que permite extraer la lectura y conservar esos datos en el sistema para su posterior análisis.

\subsection{Sprint 6: 19/11/2018 - 11/12/2018}

\subsubsection{Tareas}
\begin{itemize}
	\item 
\end{itemize}
\subsubsection{Backlog}

\subsubsection{Investigación}

\subsection{Sprint :  - }

\subsubsection{Tareas}
\begin{itemize}
	\item 
\end{itemize}
\subsubsection{Backlog}

\subsubsection{Investigación}

\section{Estudio de viabilidad}

\subsection{Viabilidad económica}

\subsection{Viabilidad legal}



\apendice{Especificación de Requisitos}

\section{Introducción}
En este apartado se presentan las funcionalidades básicas que debe cumplir el sistema desarrollado para este proyecto. En este caso se van a diferenciar dos tipos: los objetivos generales, los cuales son de obligado cumplimiento, y el resto de requisitos que se mostrarán en un catálogo debido a que  no es necesario que se cumplan todos.

\section{Objetivos generales}
Para el sistema software desarrollado en este proyecto, se fijan los siguientes objetivos generales:
\begin{itemize}
	\item Establecer una conexión con el láser.
	\item Intercambiar mensajes para obtener los datos de lectura.
	\item Implementar un analizador de dichos datos.
	\item Implementar una representación gráfica para la visualización del análisis.
\end{itemize}

\section{Catalogo de requisitos}
Para la realización de este catálogo se diferenciará entre dos tipos de requisitos: los requisitos funcionales y los requisitos no funcionales.
\subsection{Requisitos funcionales}
 \begin{itemize}
 	\item R1: Conectar con el láser.
 	\item R2: Enviar mensaje.
 	\item R3: Recoger datos.
 	\item R4: Traducir datos.
 	\item R5: Incorporar áreas al sistema.
 	\item R6: Representar gráficamente los datos del análisis.
 	\item R7: Impresión de las áreas con objetos detectados.
 \end{itemize}
\subsection{Requisitos no funcionales}
\begin{itemize}
	\item El sistema ha de detectar de la forma más precisa posible la distancia a la que se encuentra cada uno de los puntos detectados, ya se encuentren dentro o fuera de los campos de área establecidos.
	\item El software necesita estar conectado al láser para poder realizar la comunicación.
	\item Al ser una primera versión de este proyecto, es decir, al poder ser un sistema extensible y escalable, se ha codificado de forma modular.
\end{itemize}
\section{Especificación de requisitos}
\subsection{R1: Conectar con el láser.}
\textbf{Versión} 1.0\\
\textbf{Autor} Álvaro Ruifernández Palacios\\
\textbf{Descripción} El sistema tiene que conseguir crear conexión con el láser para poder interactuar con él y conseguir la información de sus lecturas.\\
\textbf{Precondición} Ha tenido que comenzar la ejecución del programa principal.\\
\textbf{Secuencia Normal} los pasos a seguir son:
\begin{enumerate}
	\item El usuario ejecuta el programa principal.
	\item Se crea el objeto de tipo Socket.
	\item Se establece la dirección IP y el número de puerto.
\end{enumerate}
\textbf{Postcondición} Al enviar el mensaje se recibirá una respuesta dependiendo de dicho mensaje\\
\textbf{Excepciones} El láser no esté físicamente conectado al PC.\\
\textbf{Importancia} Alta\\
\textbf{Comentarios} Sin comentarios\\

\subsection{R2: Enviar mensaje.}
\textbf{Versión} 1.0\\
\textbf{Autor} Álvaro Ruifernández Palacios\\
\textbf{Descripción} Se crea el mensaje para, después de haberlo codificado, enviarlo a través del socket para que este devuelva la información de lectura del entorno, es decir, se crea un mensaje de petición de datos dirigido al láser.\\
\textbf{Precondición} Se tiene que haber configurado el socket previamente (se tiene que haber cumplido R1)\\
\textbf{Secuencia Normal} Los pasos de este requisito son:
\begin{enumerate}
	\item Se crea una lista con cada uno de los caracteres que forman el mensaje codificados según corresponda (explicado en el punto 3 de la memoria).
	\item Esta lista es transformada en una array de bytes o bytearray.
	\item Este bytearray es enviado al láser a través del socket mencionado previamente.
\end{enumerate}
\textbf{Postcondición} El láser enviará la información de lectura por el mismo puerto por el que se envió la petición.\\
\textbf{Excepciones} Si el mensaje no es configurado de forma correcta (fallo en la codificación, falta de alguno de los campos...) el láser devolverá un mensaje de error.\\
\textbf{Importancia} Alta\\
\textbf{Comentarios} Sin comentarios.\\

\subsection{R3: Recoger datos.}
\textbf{Versión} 1.0\\
\textbf{Autor} Álvaro Ruifernández Palacios\\
\textbf{Descripción} Recoger y mantener en el sistema los datos recibidos de las lecturas del láser tras el paso del mensaje correcto correspondiente para su posterior análisis.\\
\textbf{Precondición} Se debe haber enviado el mensaje correcto, es decir, se debe haber cumplido pre y postcondiciones de R2.\\
\textbf{Secuencia Normal} Los pasos de este requisito son:
\begin{enumerate}
	\item El láser comienza la transmisión de información.
	\item La variable empieza a encadenar las parte de la cadena que forma los datos recibidos.
	\item  A la vez que esto, otra variable recoge la cantidad de información obtenida para recoger solo la información deseada.
	\item Cuando la variable anterior llega al máximo preestablecido por el programador, se cierra el socket quedando la primera de las variables con la información de lectura deseada.
\end{enumerate}
\textbf{Postcondición} La información de lectura que realizó el láser queda almacenada en el programa principal para posterior análisis.\\
\textbf{Excepciones} Si se interrumpe la conexión con el láser, la ejecución del programa queda en un bucle infinito al no haber completado la longitud máxima marcada por la segunda de las variables descritas en la enumeración.\\
\textbf{Importancia} Alta\\
\textbf{Comentarios} Sin comentarios.\\

\subsection{R4: Traducir datos.}
\textbf{Versión} 1.0\\
\textbf{Autor} Álvaro Ruifernández Palacios\\
\textbf{Descripción} Se utiliza la variable que almacena las lecturas del láser para extraer los datos y transformarlos en una colección de puntos para tenerlos el un formato amigable para el sistema.\\
\textbf{Precondición} Se debe haber cumplido R3, con lo que el sistema tiene almacenado la información de lectura a analizar.\\
\textbf{Secuencia Normal} Los pasos a seguir en este requisito son:
\begin{enumerate}
	\item Se divide la cadena de lectura según los caracteres de inicio y final de los mensajes.
	\item Se eliminan de la lista creada de la división todas las cadenas vacías.
	\item Se extrae a una variable la segunda de las posiciones de la lista de cadenas (ya que la primera es información de confirmación).
	\item Se elimina la información irrelevante y se divide la cadena resultante en tantos datos como pasos da el láser en cada lectura.
	\item Se asigna a cada dato un ángulo de una lista previamente configurada de la misma longitud que cantidad de datos de lectura
	\item Se traduce cada par de datos a valores cartesianos (ya que la asignación anterior creaba valores polares).
\end{enumerate}
\textbf{Postcondición} El sistema posee una lista con las coordenadas cartesianas de cada una de las lecturas.\\
\textbf{Excepciones} Si los cálculos de traducción de puntos fallan, se guardarían datos erróneos en el sistema.\\
\textbf{Importancia} Alta\\
\textbf{Comentarios} Sin comentarios.\\

\subsection{R5: Incorporar áreas al sistema.}
\textbf{Versión} 1.0\\
\textbf{Autor} Álvaro Ruifernández Palacios\\
\textbf{Descripción} Introducir en el sistema el número y las coordenadas de los límites de las áreas que el usuario desea analizar para comprobar si existen objetos.\\
\textbf{Precondición} El sistema debe estar ejecutándose sin ningún tipo de fallo.\\
\textbf{Secuencia Normal} Los pasos a seguir son:
\begin{enumerate}
	\item El sistema muestra un mensaje para que el usuario introduzca el número de áreas a introducir.
	\item Tras haber introducido el número correctamente se muestra una sucesión de mensajes para introducir cada uno de los límites de cada una de las áreas.
\end{enumerate}
\textbf{Postcondición} El sistema guarda los datos a cerca de las áreas a analizar para conocer si existen o no objetos en ellas.\\
\textbf{Excepciones} Si los datos son introducidos de manera incorrecta, serán registrados incorrectamente por lo que causará errores en el análisis.\\
\textbf{Importancia} Media\\
\textbf{Comentarios} Existe la opción de no querer introducir áreas (introduciendo 0 como número de áreas a analizar) si el usuario solo quiere observar las lecturas del láser. Esta forma de funcionamiento del programa puede ser empleada para comprobar el buen funcionamiento del sistema de traducción.\\

\subsection{R6: Representar gráficamente los datos del análisis.}
\textbf{Versión} 1.0\\
\textbf{Autor} Álvaro Ruifernández Palacios\\
\textbf{Descripción} Como bien se describe en el título, se representan gráficamente las coordenadas tanto de las lecturas traducidas como de las coordenadas de las áreas introducidas por el usuario.\\
\textbf{Precondición} Ejecución del programa principal sin ningún fallo.\\
\textbf{Secuencia Normal}Los pasos a ejecutar son:
\begin{enumerate}
	\item Se divide cada uno de los puntos (sea cual sea su origen) en dos listas, una con las coordenadas pertenecientes al eje de abscisas y la otra con las correspondientes al eje de ordenadas. A la hora de almacenar los datos separados en listas si que se tiene en cuenta su procedencia.
	\item Se introducen en la gráfica los datos de lectura del láser en color azul.
	\item Se introducen en la gráfica los datos de los límites de cada una de las áreas en color rojo.
\end{enumerate}
\textbf{Postcondición} Se muestran todos los datos a cerca de los puntos que posee el sistema en ese momento.\\
\textbf{Excepciones} Se puede producir un error a la hora de representar si los datos de las áreas se encuentran mal introducidos.\\
\textbf{Importancia} Media\\
\textbf{Comentarios} Sin comentarios\\

\subsection{R7: Impresión de las áreas con objetos detectados.}
\textbf{Versión} 1.0\\
\textbf{Autor} Álvaro Ruifernández Palacios\\
\textbf{Descripción} Se imprimen los mensajes de las áreas, es decir, el sistema dice si se ha encontrado algún punto dentro de estas áreas, lo que significaría que hay un objeto en alguna de ellas.\\
\textbf{Precondición} Los datos de las áreas deben de estar correctamente introducidos y no haber aparecido ningún mensaje de error que hubiese detenido la ejecución.\\
\textbf{Secuencia Normal}Los pasos a seguir son:
\begin{enumerate}
	\item Se comprueba si las coordenadas de alguno de los puntos de lectura se encuentran entre los límites de alguna de las áreas.
	\item En caso positivo aparece un mensaje con la posición de ese punto.
	\item En caso negativo aparece un mensaje diciendo que no se encontró ningún objeto.
\end{enumerate}
\textbf{Postcondición} El usuario es informado de si en las áreas que deseaba analizar se ha encontrado o no algún objeto, indicando la posición del objeto encontrado en caso de que así fuese.\\
\textbf{Excepciones} Se pueden producir errores dependiendo de la corrección con la que se hayan introducido las coordenadas de las áreas.\\
\textbf{Importancia }Alta\\
\textbf{Comentarios} Sin comentarios.\\


\apendice{Especificación de diseño}

\section{Introducción}

En esta sección se describirá el proceso de diseño de cada uno de los aspectos del proyecto que forman la base para conseguir alcanzar los objetivos propuestos para este proyecto.\\

\section{Diseño de datos}

En este proyecto se emplea un único tipo de datos, las coordenadas cartesianas, las cuales se dividen en dos grupos dependiendo de su procedencia:
\begin{itemize}
	\item Puntos de lectura: grupo formado por las lecturas que provienen de las lecturas del láser y cuya traducción de formato cadena (formato de llegada al sistema) a coordenadas cartesianas ya ha sido descrita en apartados anteriores.
	\item Límites de áreas: grupo formado por los límites de las áreas introducidas por el usuario, las cuales no necesitan traducción ya que son introducidas directamente como coordenadas cartesianas.
\end{itemize}
Ambos grupos se tratan de la misma forma dentro del sistema, ya que se crean listas de instancias de la clase \textit{Punto}, creada en este proyecto.

\section{Diseño procedimental}

Para este apartado se ha desarrollado el siguiente diagrama de secuencia, en el que se describe el proceso de ejecución del sistema.Como se puede observar en la imágen del diagrama, la única interacción que tiene el usuario es la introducción del número de áreas y las coordenadas de cada uno de sus límites.
\imagenPagCompleta{diagramaSec}{Diagrama de secuencia del sistema}

\section{Diseño arquitectónico}

En cuanto a la arquitectura del proyecto se destaca su simplicidad. Estos es debido a que la base del proyecto es la creación de un proceso, por lo que no posee clases como interfaces u otros elementos a dividir en varias clases.
\imagen{diagramaClases}{Diagrama de clases del sistema}
\apendice{Documentación técnica de programación}

\section{Introducción}
En este apartado se mostrarán todas las herramientas que necesita un programador para poder utilizar y trabajar en este proyecto. Como primer paso, antes de cualquier herramienta, se necesita acceso al proyecto.\\
\\
Este proyecto es accesible a través de la URL: \url{https://github.com/arp0054/TFG_ObjDetectorViaLaser}\\

\section{Estructura de directorios}
Debido a las diferentes formas de implementación que este proyecto ha tenido a lo largo de su desarrollo, se pueden observar una gran variedad de directorios:
\begin{itemize}
	\item DataSheets: en este directorio se encuentran los documentos técnicos y el protocolo de comunicación utilizado por el láser, los cuales han servido para obtener una gradn cantidad de datos que han ayudado al desarrollo del proyecto (dirección IP y número de puerto por defecto, formas de encriptación y desencriptación de los datos...).
	\item Diagramas: en este caso el nombre es muy descriptivo, en este directorio se encuentran los diagramas de clases y de secuencia, tanto los iniciales, desarrollados para las primeras fases del proyecto, como los finales. Estos diagramas se distinguen por estar separados en los subdirectorios "Iniciales" y "Finales".
	\item Dispositivo: también perteneciente a las fases iniciales, en él aparecen el manual de usuario del MTX-GTW, con detalladas explicaciones a cerca de sus necesidades eléctricas así como todas las aplicaciones que el hardware de este dispositivo es capaz de realizar.
	\item Instalación: en este caso, aparecen tanto imágenes como documentos a cerca de como instalar las librerías las cuales permitían la ejecución remota en el sistema descrito en el punto anterior así como otras utilidades necesarias para el desarrollo del proyecto en su primera versión (útiles para su uso con C$++$)
	\item Memoria: como su mismo nombre indica, este es el directorio en el que se guardan todos los archivos e imágenes que conforman la memoria de este proyecto, así como su sección de anexos.
	\item Python: en este directorio se alamcenan todos los archivos del sistema codificados en el lenguaje de programación con el mismo nombre, los cuales forman la versión actual y funcional del proyecto. Hay dos archivos con el mismo nombre " objectDetectorViaLaser" aunque con distinta extensión, dependiendo si es un archivo para usar con Jupyter Notebook (extensión ".ipynb"), el cual se utilizó primero para desarrollo y después solo para alguna prueba, o para utilización en Spyder(extensión ".py"). De la utilización de estos archivos aparece la generación automática de los dos subdirectorios pertenecientes a esta parte del repositorio.
	\item TFG\_ 1718: de nuevo un directorio de las primeras versiones, ya que este alberga todas las clases, ya sean funcionales o librerias, que fueron usadas en las primeras versiones para codificar el sistema objetivo en el lenguaje de programación C$++$ .
	\item Drivers: este directorio almacena simplemente el driver utilizado para la conexión con el láser vía USB.
\end{itemize}
Como se ha podido observar,algunos de los directorios pertenecen a las primeras versiones. Esto viene dado por la motivación ta expuesta en el punto 7 de la memoria, donde una de las 
\section{Manual del programador}
En esta sección se explica como se configura el entorno de trabajo para este proyecto.
\subsection{Python}
Como primer paso para esta configuración se debe descargar el lenguaje de programación que se va a emplear, en este caso Python. Para este proyecto no se especifica ninguna versión en especial pero para su desarrollo se ha empleado la versión número 3.6.3.\\
Para  su descarga visitar: \url{https://www.python.org/downloads/release/python-363/}\\
\\
Después de descargar e instalar este lenguaje, el cual tiene incorporado el manejador de paquetes \textit{pip}, se descargan los módulos o paquetes que se emplean sobre todo en la representación gráfica con los siguientes comandos:
\begin{itemize}
	\item pip install matplotlib.
	\item pip install numpy
\end{itemize}
Estas bibliotecas son las que permiten preparar y representar los datos a mostrar de forma gráfica tras haber pasado por el proceso de análisis y traducción.\\

\subsection{IDE}
Para poder modificar o crear nuevo código en este proyecto se necesita un entorno de desarrollo integrado o IDE. Aunque Python ya tiene el suyo propio, se ha optado por el uso de otros debido a que son más amigables y algunos de ellos se han empleado a lo largo del grado. Estos entornos son: Jupyter Notebook y Spyder.\\
\\
Para descargar ambos entornos se debe instalar la aplicación de Anaconda Navigator, la cual incorpora a ambos. Para descargarlo hay que visitar la página web: \url{https://anaconda.org/anaconda/anaconda-navigator} \\
\\
La principal diferencia entre ambos es que Jupyter utiliza una interfaz web, mientras que Spyder posee su propia interfaz.\\
\imagen{interfazJupyter}{Interfaz gráfica de Jupyter}
\imagen{interfazSpyder}{Interfaz gráfica de Spyder}

\subsection{Git}
Este es el sistema escogido para el control de versiones del proyecto. Al no estar incluido por defecto en Windows, se debe instalar para poder utilizarlo, al igual que el resto de programas utilizados. Para poder descargarlo se ha de acceder a la pagina cuya URL es: \url{https://git-scm.com/}\\
Dependiendo de las preferencias del programador se pueden instalar diferentes formas de instalación (consola de bash, de Windows..).

\subsection{GitKraken}
Este cliente de Git es el escogido para el mantenimiento de versiones en el repositorio, es decir, el cliente con el que se subirán cada uno de los cambios realizados al repositorio y con el  que se descargarán las últimas versiones del proyecto desde cualquier dispositivo empleado para el desarrollo de cualquiera de las partes del proyecto. Para descargarlo se debe acceder a: \url{https://www.gitkraken.com/download}\\
\\
Una vez descargado, hemos de acceder a GitHub y buscar el botón \textit{Clone or download} y copiar la URL que aparece tras presionarlo (\url{https://github.com/arp0054/TFG_ObjDetectorViaLaser.git}). Una vez copiada, se realizan los siguientes pasos en GitKraken: $boton con el icono de directorio \rightarrow CLONE \rightarrow Clone with URL \rightarrow se selecciona el directorio donde se van a guardar los archivos del directorio en el PC \rightarrow se pega la URL previamente copiada \rightarrow Clone repo! $\\
\\
Una vez completados esos pasos se habrá descargado el repositorio por lo que se puede usar como cualquier otro proyecto.\\
\imagen{interfazGitKraken}{Interfaz gráfica de GitKraken}

\section{Compilación, instalación y ejecución del proyecto}
Si se han cumplimentado todos los pasos anteriores en cuanto a instalaciones del entorno (Python y sus librerías) el sistema es capaz de compilar y ejecutar el software de este proceso (ya que no tiene proceso de instalación como tal).\\
\\
Para este proceso de ejecución se puede realizar de tres formas, atendiendo a los 3 entornos posibles:
\begin{enumerate}
	\item por consola de comandos (CMD en Windows): para este proceso se ha de abrir el programa \textit{Símbolo del sistema} (instalado por defecto en el sistema operativo), utilizar el comando \textit{cd} para colocarnos en la carpeta donde se encuentran todas las clases de Python y escribir el comando \textbf{objectDetectorViaLaser.py}
	\item Por Jupyter Notebook: se debe buscar el archivo objectDetectorViaLaser.ípynb desde la interfaz de Jupyter, la cual es similar a un explorador de archivos, y de ahí presionar el botón de ejecución tantas veces como bloques de código posea el archivo.
	\item Por Spyder: se debe de buscar el archivo desde el explorador de archivos, similar a cualquier otro programa, y ejecutarlo una única vez ya que en este caso no existe separación. 
\end{enumerate}
Ya que son equipotenciales, se deja a decisión del programador la forma en la que se ejecuta el proceso, ya que l compilación es automática en todos ellos.


\section{Pruebas del sistema}
A lo largo del desarrollo de este proyecto se han ido desarrollando una serie de pruebas las cuales, según su resultado, ayudaban al desarrollo de las diferentes fases o a la detección de errores. En este apartado se expondrán las pruebas realizadas y los resultados obtenidos en las mismas. Algunas de las pruebas descritas en este apartado (3 - 6) no han sido únicas, es decir, se han realizado varias veces, pero, debido a que los datos son idénticos en cada uno de los apartados, se ha decidido exponer como caso de prueba.\\
\\
Además de estas pruebas se probó el funcionamientos del sistema con el simulador, lo que con las primeras pruebas hizo que se observase y se corrigiese un fallo en el manejo de los hilos de cada una de las ejecuciones del software principal, el cual se corrigió y en la versión actual no presenta ningún problema.
\subsection{Prueba 1: Comunicación con MTX-GTW}
\textbf{Descripción} Se prueba a utilizar el proyecto de C$++$ en el cual solo se posee un archivo que imprime por pantalla "Hola Mundo!" para que, con la configuración del proyecto se ejecute en el MTX-GTW y verificar la comunicación\\
\textbf{Entorno} El ordenador se encuentra conectado al dispositivo antes mencionado vía cable de red y este a su vez se encuentra conectado al cargador, creado por el alumno, para poder funcionar.\\
\textbf{Resultado esperado} Impresión por pantalla del proceso de ejecución remota (llamada a la dirección IP, mensaje de confirmación...) y la impresión por pantalla del mensaje presente en el proyecto.\\
\textbf{Resultado obtenido} Se obtienen todos los mensajes esperados.\\
\textbf{Conclusiones} Se observa que por el momento se aprecia una configuración y comunicación con el MTX-GTW correcta.
\textbf{Acciones derivadas}  proseguir con el proyecto.\\
\subsection{Prueba 2: ejecución en MTX-GTW}
\textbf{Descripción}Después de un tiempo desarrollando partes del código, se decide ejecutar de forma remota dicho código para comprobar el buen funcionamiento del mismo. Al ser clases muy sencillas, se decide crear impresiones de pantalla para poder visualizarlo \\
\textbf{Entorno} mismo entono que Prueba 1.\\
\textbf{Resultado esperado} Impresión de mensajes de comunicación con el dispositivo conectado e impresiones de todos los mensajes.\\
\textbf{Resultado obtenido} Múltiples mensajes de error indicando fallos en la comunicación.\\
\textbf{Conclusiones} Se observan fallos en la comunicación con el dispositivo y en la configuración del proyecto.\\
\textbf{Acciones derivadas}  Se comienza un análisis de la configuración del proyecto, con lo que se descubre que faltan algunas librerías las cuales no se pueden encontrar debido a que se encuentran desactualizadas y no se encuentra forma de acceder a ellas.\\

\subsection{Prueba 3: Comunicación con RealTerm}
\textbf{Descripción} Una vez se ha eliminado al idea de la ejecución remota se comienzan las pruebas de comunicación con el láser, estas primeras pruebas son a cerca de la creación de los mensajes que se emplearían más tarde en el código. Para ello, se envía al láser el mensaje "2,0,0,0,E,A,R,0,2,0x83,0x00,3" ya que es la estructura según establecida para el mensaje de comunicación contínua, que es es que se va a querer usar para el programa principal.Este mensaje es enviado como números.\\
\textbf{Entorno} Se conectan láser y PC vía cable Ethernet. A su vez, el láser está conectado al cargador, creado por el alumno, para poder usarlo en la comunicación. En este caso, los alrededores del láser no son relevantes debido a que solo se está comprobando el intercambio de información\\
\textbf{Resultado esperado} El láser devuelve una sucesión continua o muy extensa de los datos de lectura.\\
\textbf{Resultado obtenido} No se recibe ningún mensaje o respuesta y si se reciben son mensajes de error\\
\textbf{Conclusiones} La codificación del mensaje es incorrecta, por lo que necesita revisión.\\
\textbf{Acciones derivadas}  Se revisa la codificación y se observa las dos diferentes codificaciones en el mismo mensaje. Se estudia como crear ese mensaje.\\

\subsection{Prueba 4: Segunda comunicación con RealTerm}
\textbf{Descripción} Tras analizar las codificaciones y la forma en la que se deben transmitir cada una de las codificaciones necesarias para la creación del mensaje. en este caso el mensaje creado es "2,48,48,48,69,65,82,48,50,0x00,0x83,3" y en este caso se transmiirá como caracateres ASCII. En cuanto al entorno, como en la anterior prueba, es irrelevante debido a que lo que se está comprobando es la forma de comunicación, no la información en sí.\\
\textbf{Entorno} El mismo que en la prueba 3.\\
\textbf{Resultado esperado} Conjunto grande de datos formados por varias lecturas del entorno.\\
\textbf{Resultado obtenido} Conjunto de datos esperado, aunque con fallo por llenado de buffer de entrada configurado por RealTerm\\
\textbf{Conclusiones} El mensaje ha sido creado correctamente\\
\textbf{Acciones derivadas}  Se desarrolla el código en Python para la recepción, análisis de datos y representación de los mismos.

\subsection{Prueba 5: Traducción y representación de datos}
\textbf{Descripción} Una vez creado el código del programa principal en Python se desea representar los datos correspondientes a la lectura de datos con el láser apuntando hacia una pared. \\
\textbf{Entorno} En este caso se tiene en cuenta el entorno debido a que se va a representar la lectura de una línea recta correspondiente a la pared. En cuanto al equipo, es idéntico al de las 2 pruebas anteriores.\\
\textbf{Resultado esperado} Puntos dispersos en los laterales pero una línea de puntos al frente, representando la pared anteriormente mencionada.\\
\textbf{Resultado obtenido} Puntos dispersos, creando círculos concéntricos, los cuales se extendian por los 360º cuando los lidar solo son capaces de leer 270º\\
\textbf{Conclusiones} Fallo en la traducción o modificación de los datos recibidos por las lectura del láser pero, por los resultados obtenidos, se deduce que el problema es de la traducción de coordenadas.\\
\textbf{Acciones derivadas}  Se analiza el código en busca de dicho error y, como ya se había deducido en las conclusiones, el problema surgía del paso de coordenadas polares a cartesianas debido a una confusión con los datos, lo cual se rectificó.

\subsection{Prueba 6: Segunda prueba de Traducción y representación de datos}
\textbf{Descripción} Una vez analizadas y corregidas todos los métodos de traducción de puntos utilizados, se utiliza la caja del propio láser sobre él mismo para que el sistema lo represente de forma gráfica.\\
\textbf{Entorno} El entorno es similar al que se expone en la prueba anterior, salvo el aspecto de la sustitución de la pared por la caja para representarla gráficamente.\\
\textbf{Resultado esperado} Dibujo de la silueta de la caja en forma de agrupación de puntos.\\
\textbf{Resultado obtenido} Tras ejecutar el programa principal se obtiene la siguiente imagen:\\
\imagen{pruebaCaja}{Prueba de la representación de la caja como lectura de entorno}
\textbf{Conclusiones} La traducción de datos funciona de forma correcta.\\
\textbf{Acciones derivadas}  Se realizaron varias repeticiones de esta prueba para la obtención de varias lecturas, las cuales se han empleado para la creación de la lista de lecturas manejada por el servidor para la simulación de comportamiento que conforma la base de su funcionamiento. Aqui un ejemplo
\imagenGrande{prueba6E}{Entorno de una de una prueba correcta}
\imagenGrande{prueba6G}{Gráfica resultado de una prueba correcta}
\apendice{Documentación de usuario}

\section{Introducción}
En este apartado se verán todos los aspectos necesarios para poder instalar y usar el software de forma correcta.

\section{Requisitos de usuarios}
Para que un usuario pueda usar este sistema, debe poseer antes el hardware necesario para ello:
\begin{itemize}
	\item Un ordenador: en el cual se va a ejecutar el software. Este necesita tener como sistema operativo Windows 10, ya sea como sistema operativo principal o el un sistema virtualizado, aunque se recomienda más la primera de las opciones debido a que se evitan posibles fallos a causa de los servicios ofrecidos por las máquinas virtuales.
	\item Un láser: con el que se van a realizar las lecturas de entorno. En este caso se recomienda el mismo utilizado para el desarrollo (Hokuyo Safety Laser Scanner  o UAM-05LP-T301) ya que el programa ha sido diseñado y se encuentra configurado para él. Si este láser se cambiase por otro se tendrían que cambiar algunos de los parámetros más importantes del sistema (mensajes utilizados, forma de traducción, configuración del socket...).
	\item Cable Ethernet: empleado para la comunicación de los dos elementos anteriormente nombrados. En este caso, la elección del cable se deja a elección del usuario, aunque se recomienda el creado por la empresa que desarrolló del láser o uno con características similares.
	\item Transformador eléctrico: que suministre 24 V de corriente, voltaje requerido por el láser para poder funcionar correctamente. Para esto se utilizará también un par de empalmes de cables para conectar los cables del láser (los cuales no poseen cabezal convencional) a los cables del transformador.
\end{itemize}
\section{Instalación}
Para este paso se deben de seguir los pasos descritos en el \textit{Apéndice D} a cerca de la instalación de Python y sus librerías. Una vez hecho esto, el único paso que se debe dar es:
\begin{itemize}
	\item Se accede al repositorio con el enlace presente también en el apéndice antes mencionado en formato ZIP.
	\item Se descomprime el proyecto completo o el directorio \textit{Python} (a elección del usuario).
	\item El usuario ya tiene el proyecto listo para ser operativo
\end{itemize}

\section{Manual del usuario}
Para poder utilizar de forma correcta este software, una vez completado el proceso de instalación, es ejecutar el código tanto del servidor (si se desea o no se tiene acceso a un láser) como el del programa principal (en este orden para su correcto funcionamiento en el caso de no tener el láser conectado físicamente) en consola de comando o, si el usuario lo tiene instalado, un IDE que permita la ejecución de Python. En este apartado se va a explicar paso por paso la realización de este proceso (en este caso se expondrá la forma de uso desde consola de comandos).\\
\\
Como primer paso, se deben abrir dos consolas de comandos con la orden \textit{cmd}. Una vez abiertas se abre el directorio donde se encuentre el proyecto, y más concretamente el directorio con los archivos Python.\\
\imagen{locDir}{colocación de consolas en directorio Python}
En la primera de las consolas es donde se ejecutará el servidor que actuará como simulador del láser (en caso de que el láser se encuentre conectado este paso se puede omitir) con la ejecución del comando \textit{py Simulator.py}. Por otra parte, en la segunda consola se realizará la ejecución del programa principal a través de la orden \textit{py objectDetectorViaLaser.py}.\\
\imagen{exeProc}{Ejecución de comandos en ambas consolas}
Tras estas ejecuciones, el programa principal muestra por pantalla un mensaje para que el usuario introduzca el número de áreas que desea. Después de  que el usuario haya completado esta tarea, se le pedirá la introducción de las coordenadas de cada una de las áreas con un mensaje indicativo de la unidad de medida en la que se van a interpretar los datos (milímetros) y la orden de separa los datos utilizando la coma (,).\\
\imagen{intrArea}{Introducción de áreas}
Una ve introducidas todas las áreas el sistema se intentará conectar al láser utilizando su dirección IP y su número de puerto. Al no haber podido por no estar conectado físicamente al láser, el sistema se conecta a la dirección y puerto del servidor (en caso de estar conectado al láser la conexión se realiza con este). Tras realizar esta conexión se ejecuta el resto del programa principal, haciendo aparecer el mensaje correspondiente de cada una de las áreas y la información correspondiente a la ventana que emerge a continuación con la representación gráfica de los datos, tanto de lectura simulada (o real en caso de conexión con el láser) y las áreas insertadas por el usuario.\\
\imagen{pingYEx}{Visualización de datos de ejecución}
\imagen{repGr}{Gráfica de representación de datos}
Para poder mostrar de forma más visual este proceso, tanto de la descarga del proyecto como de la ejecución del mismo, se puede ver el siguiente vídeo: \url{https://github.com/arp0054/TFG_ObjDetectorViaLaser/tree/master/Videos}


\bibliography{bibliografiaAnexos}
\bibliographystyle{plain}

\end{document}
